% Part: first-order-logic
% Chapter: syntax-and-semantics
% Section: Semantic Notions

\documentclass[../../../include/open-logic-section]{subfiles}

\begin{document}

\olfileid[fr]{fol}{syn}{sem}
\olsection{Notions sémantiques}

\begin{explain}
La définition des !!{structure}s pour les langages du premier ordre
permet de formuler des propriétés sémantiques fondamentales des
énoncés et des relations entre eux. La plus simple est la
\emph{validité} d'un énoncé : il est valide s'il est satisfait dans
toute !!{structure}. Sa satisfaction ne dépend donc pas de la façon
dont ses symboles non logiques sont interprétés. Les énoncés valides
sont aussi appelés \emph{vérités logiques} : ils sont vrais,
c'est-à-dire satisfaits, dans toute !!{structure}. Leur vérité
ne dépend que des symboles logiques qu'ils contiennent et de leur
structure syntaxique, et non des symboles non logiques ni de leur
interprétation.
\end{explain}

\begin{defn}[Validité]
Un énoncé~$!A$ est \emph{valide}, ce que nous notons $\Entails !A$,
si et seulement si $\Sat{M}{!A}$ pour toute
!!{structure}~$\Struct M$.
\end{defn}

\begin{defn}[Conséquence logique]
Un ensemble d'énoncés~$\Gamma$ \emph{entraîne} un énoncé~$!A$,
ce que nous notons $\Gamma \Entails !A$, si et seulement si, pour
toute !!{structure}~$\Struct M$ telle que $\Sat{M}{\Gamma}$,
on a $\Sat{M}{!A}$.
\end{defn}


\begin{defn}[Satisfaisabilité]
Un ensemble d'énoncés~$\Gamma$ est \emph{satisfaisable} s'il existe
!!a{structure}~$\Struct M$ telle que $\Sat{M}{\Gamma}$. S'il n'est
pas satisfaisable, il est dit \emph{insatisfaisable}.
\end{defn}

\begin{prop}
Un énoncé~$!A$ est valide si et seulement si
$\Gamma \Entails !A$ pour tout ensemble d'énoncés~$\Gamma$.
\end{prop}

\begin{proof}
Dans le sens direct, supposons $!A$ valide et soit $\Gamma$ un
ensemble d'énoncés. Si une !!{structure}~$\Struct M$ satisfait
$\Gamma$, alors $\Sat{M}{!A}$ par validité de~$!A$; donc
$\Gamma \Entails !A$.

Pour la contraposée du sens réciproque, supposons $!A$ non valide.
Il existe alors !!a{structure}~$\Struct M$ telle que
$\Sat/{M}{!A}$. Prenons $\Gamma = \{ \ltrue \}$. Comme
$\ltrue$ est valide, $\Sat{M}{\Gamma}$. Cette !!{structure}
satisfait donc~$\Gamma$ sans satisfaire~$!A$ :
$\Gamma$ n'entraîne pas~$!A$.
\end{proof}

\begin{prop}
\ollabel{prop:entails-unsat}
$\Gamma \Entails !A$ si et seulement si
$\Gamma \cup \{\lnot !A\}$ est insatisfaisable.
\end{prop}

\begin{proof}
Dans le sens direct, supposons $\Gamma \Entails !A$ et, par
l'absurde, qu'il existe !!a{structure}~$\Struct M$ telle que
$\Sat{M}{\Gamma \cup \{ \lnot !A \}}$. Puisque
$\Sat{M}{\Gamma}$ et $\Gamma \Entails !A$, nous avons
$\Sat{M}{!A}$. Mais $\Sat{M}{\Gamma\cup \{ \lnot !A \}}$ donne
aussi $\Sat{M}{\lnot !A}$, donc $\Sat/{M}{!A}$ : contradiction.
Une telle !!{structure}~$\Struct M$ n'existe pas; l'ensemble
$\Gamma \cup \{ \lnot !A \}$ est insatisfaisable.

Réciproquement, supposons $\Gamma \cup \{ \lnot !A \}$
insatisfaisable. Pour toute !!{structure}~$\Struct M$, on a alors
$\Sat/{M}{\Gamma}$ ou $\Sat{M}{!A}$. Ainsi, toute !!{structure}
$\Struct M$ qui satisfait~$\Gamma$ satisfait aussi~$!A$;
donc $\Gamma \Entails !A$.
\end{proof}

\begin{prob}
\begin{enumerate}
\item Montrer que $\Gamma \Entails \bot$ si et seulement si
  $\Gamma$ est insatisfaisable.
\item Montrer que $\Gamma \cup \{!A\} \Entails \bot$ si et seulement si
  $\Gamma \Entails \lnot !A$.
\item Supposons que $!A$ n'ait pas de variable libre autre que $x$ et que $c$ soit une constante n'apparaissant ni dans~$!A$ ni dans~$\Gamma$.\footnote{Précision éditoriale : la conséquence logique a été définie ici pour des énoncés. Cette condition garantit que la quantification et le remplacement par la constante produisent bien des énoncés.}
  Montrer que $\Gamma \Entails \lforall[x][!A]$ si et seulement si
  $\Gamma \Entails
  \Subst{!A}{c}{x}$.
\end{enumerate}
\end{prob}

\begin{prop}
Si $\Gamma \subseteq \Gamma'$ et $\Gamma \Entails !A$, alors
$\Gamma' \Entails !A$.
\end{prop}

\begin{proof}
Supposons $\Gamma \subseteq \Gamma'$ et $\Gamma \Entails !A$.
Soit $\Struct M$ une structure telle que $\Sat{M}{\Gamma'}$.
Alors $\Sat{M}{\Gamma}$ et, puisque $\Gamma \Entails !A$,
on obtient $\Sat{M}{!A}$. Toute structure satisfaisant~$\Gamma'$
satisfait donc~$!A$, ce qui donne $\Gamma' \Entails !A$.
\end{proof}


\begin{thm}[Théorème sémantique de déduction]
\ollabel{thm:sem-deduction}
$\Gamma \cup \{!A\} \Entails !B$ si et seulement si
$\Gamma \Entails !A \lif !B$.
\end{thm}

\begin{proof}
Dans le sens direct, supposons
$\Gamma \cup \{ !A \} \Entails !B$ et soit $\Struct M$ une
!!{structure} telle que $\Sat{M}{\Gamma}$. Si $\Sat{M}{!A}$,
alors $\Sat{M}{\Gamma \cup \{ !A \} }$; l'hypothèse donne donc
$\Sat{M}{!B}$. Par conséquent, $\Sat{M}{!A \lif !B}$ et
$\Gamma \Entails !A \lif !B$.

Réciproquement, supposons $\Gamma \Entails !A \lif !B$ et soit
$\Struct M$ une !!{structure} telle que
$\Sat{M}{\Gamma \cup \{ !A \}}$. Alors $\Sat{M}{\Gamma}$,
donc $\Sat{M}{!A \lif !B}$. Puisque $\Sat{M}{!A}$ également,
on a $\Sat{M}{!B}$. Toute structure satisfaisant
$\Gamma \cup \{ !A \}$ satisfait donc~$!B$ :
$\Gamma \cup \{ !A \} \Entails !B$.
\end{proof}

\begin{prop}\ollabel{prop:quant-terms}
  Soient $\Struct{M}$ !!a{structure}, $!A(x)$ !!a{formula} dont
  la seule variable libre est~$x$, et $t$~un terme clos. Alors :
  \begin{tagenumerate}{prvEx,prvAll}
    \tagitem{prvEx}{$!A(t) \Entails \lexists[x][!A(x)]$}{}
    \tagitem{prvAll}{$\lforall[x][!A(x)] \Entails !A(t)$}{}
  \end{tagenumerate}
\end{prop}

\begin{proof}
  \begin{tagenumerate}{prvEx,prvAll}
    \tagitem{prvEx}{%
      \iftag{probEx}{Exercice.}{Supposons $\Sat{M}{!A(t)}$.
        Soit $s$ une assignation de variables telle que
        $s(x) = \Value{t}{M}$. Comme $!A(t)$ est !!a{sentence},
        $\Sat{M}{!A(t)}[s]$. D'après la
        \olref[ext]{prop:ext-formulas}, $\Sat{M}{!A(x)}[s]$; puis,
        d'après la
        \olref[ass]{prop:sat-quant},
        $\Sat{M}{\lexists[x][!A(x)]}$.}}{}
    \tagitem{prvAll}{%
      \iftag{probAll}{Exercice.}{Supposons
        $\Sat{M}{\lforall[x][!A(x)]}$. Soit $s$ une assignation
        de variables telle que $s(x) = \Value{t}{M}$. D'après la
        \olref[ass]{prop:sat-quant}, $\Sat{M}{!A(x)}[s]$; puis,
        d'après la \olref[ext]{prop:ext-formulas},
        $\Sat{M}{!A(t)}[s]$. Puisque $!A(t)$ est !!a{sentence},
        \olref[ass]{prop:sentence-sat-true} donne
        $\Sat{M}{!A(t)}$.}}{}
  \end{tagenumerate}
\end{proof}

\begin{probtag}{probEx,probAll}
  Compléter la démonstration de la
  \olref[fol][syn][sem]{prop:quant-terms}.
\end{probtag} 

\end{document}
